\chapter{PHÂN TÍCH Ý NGHĨA NGHIỆP VỤ \& ĐỀ XUẤT CHIẾN LƯỢC QUẢN TRỊ}

\section{Phân tầng Rủi ro Tín dụng \& Ma trận Duyệt vay Tự động}
Dựa trên xác suất vỡ nợ $p = P(\text{Default}=1 | X)$ do mô hình Gold Layer trả về, hệ thống phân chia hồ sơ thành \textbf{4 Phân tầng Rủi ro (Risk Segments)} theo chuẩn mực quản trị rủi ro Basel \cite{basel2017} và tự động hóa quy trình ra quyết định:

\begin{table}[h!]
\centering
\caption{Ma trận Quyết định Phê duyệt Tín dụng Tự động}
\label{tab:decision_matrix}
\begin{tabularx}{\textwidth}{|l|c|X|c|}
\hline
\textbf{Phân tầng Rủi ro} & \textbf{Ngưỡng Xác suất ($p$)} & \textbf{Quyết định Duyệt Vay} & \textbf{Lãi suất đề xuất} \\ \hline
\textbf{1. Rủi ro Thấp} & $p < 15\%$ & \textbf{Duyệt Tự Động 100\% (Auto-Approve)} & 8.5\% / năm \\ \hline
\textbf{2. Rủi ro Trung Bình} & $15\% \le p < 35\%$ & \textbf{Duyệt Hạn Mức Tự Động 70\%} & 11.5\% / năm \\ \hline
\textbf{3. Rủi ro Cao} & $35\% \le p < 60\%$ & \textbf{Thẩm Định Thủ Công + Bảo Đảm} & 14.5\% / năm \\ \hline
\textbf{4. Rủi ro Rất Cao} & $p \ge 60\%$ & \textbf{Từ Chối Cho Vay (Auto-Reject)} & N/A \\ \hline
\end{tabularx}
\end{table}

% SƠ ĐỒ QUY TRÌNH PHÊ DUYỆT TÍN DỤNG TỰ ĐỘNG
\begin{figure}[h!]
\centering
% [GHI CHÚ DÀNH CHO NGHƯỜI DÙNG: CHÈN SƠ ĐỒ QUY TRÌNH DUYỆT VAY VÀO ĐÂY BẰNG LỆNH \includegraphics[width=0.85\textwidth]{so_do_quy_trinh_duyet_vay.png}]
\framebox[\textwidth][c]{\vspace{2.5cm} \textbf{[SƠ ĐỒ 5.1: QUY TRÌNH PHÊ DUYỆT TÍN DỤNG TỰ ĐỘNG VỚI MA TRẬN 4 PHÂN TẦNG]} \vspace{2.5cm}}
\caption{Sơ đồ Quy trình Phê duyệt Tín dụng Tự động dựa trên xác suất dự báo $p$}
\label{fig:sodo_decision_process}
\end{figure}

\section{Mô hình Định giá Lãi suất theo Rủi ro (Risk-Based Pricing)}
Áp dụng nguyên lý bù đắp rủi ro tín dụng (Risk Premium) \cite{arminger1997}:
$$\text{Lãi suất vay đề xuất (\%)} = \text{Lãi suất cơ bản (8.5\%)} + (p \times 15\%)$$

Khách hàng có xác suất nợ xấu $p = 5\%$ sẽ nhận mức lãi suất $9.25\%$/năm, trong khi khách hàng có $p = 25\%$ sẽ chịu mức lãi suất $12.25\%$/năm để bù đắp chi phí rủi ro cho ngân hàng.

\section{Đánh giá tác động tài chính \& Lợi nhuận đầu tư (ROI)}
\begin{enumerate}
    \item \textbf{Giảm thiểu tỷ lệ NPL:} loại bỏ hơn $80\%$ hồ sơ vỡ nợ tiềm ẩn ở nhóm Rủi ro Rất Cao, giảm tỷ lệ NPL tổng thể xuống dưới \textbf{$4.2\%$}.
    \item \textbf{Tiết kiệm chi phí trích lập dự phòng:} giảm hàng chục tỷ đồng tiền trích lập dự phòng rủi ro hàng năm theo tiêu chuẩn NHNN và Basel \cite{basel2017}.
    \item \textbf{Tăng tốc độ xử lý:} phê duyệt tự động cho $>65\%$ số lượng hồ sơ trong \textbf{<5 phút}.
\end{enumerate}

\section{Khuyến nghị quản trị cho Ban Giám đốc}
\begin{itemize}
    \item Đơn giản hóa quy trình giải ngân cho Nhóm 1 để mở rộng thị phần.
    \item Siết chặt kiểm soát chứng từ đối với các mục đích vay Medical và Debt Consolidation.
    \item Tích hợp thêm dữ liệu hành vi (Alternative Data) như hóa đơn điện nước và mua sắm TMĐT.
\end{itemize}
\newpage
